\documentclass[12pt,a4paper]{article}
\usepackage[utf8]{inputenc}
\usepackage[T2A]{fontenc}
\usepackage[russian]{babel}
\usepackage{amsmath,amssymb,amsthm}
\usepackage{booktabs,array}
\usepackage{geometry}
\usepackage{microtype}
\geometry{margin=2.3cm}

\newtheorem{proposition}{Предложение}
\newtheorem{protocol}{Протокол}
\newtheorem{nogocriterion}{Критерий провала}
\DeclareMathOperator{\Tr}{Tr}

\title{SU(5)-действие и относительный determinant:\\
два последовательных gate-теста}
\author{}
\date{4 августа 2026 г.}

\begin{document}
\maketitle

\section{Проверяемая задача}

Предыдущий rank-selector обнаружил красивое соответствие между коэффициентами
атласных формул и разложением
\begin{equation}
 \mathbf{24}=({\bf8},{\bf1})_0\oplus({\bf1},{\bf3})_0
 \oplus({\bf1},{\bf1})_0\oplus({\bf3},{\bf2})_{-5/6}
 \oplus(\overline{\bf3},{\bf2})_{+5/6}.
\end{equation}
Настоящий gate отвечает на два более строгих вопроса:
\begin{enumerate}
 \item выводятся ли rank-формулы из одного канонического $SU(5)$-действия;
 \item может ли scheme-safe twisted/untwisted determinant исправить gauge-сектор.
\end{enumerate}

\section{Gate I: канонический gauge trace}

В фундаментальном представлении
\begin{equation}
 \Tr(T_3^2)=\Tr(T_2^2)=\frac12,
 \qquad
 \Tr(Y^2)=\frac56.
\end{equation}
Поэтому один $SU(5)$-инвариантный след фиксирует
\begin{equation}
 k_Y=\frac53,
 \qquad
 \sin^2\theta_W=\frac{1}{1+k_Y}=\frac38
\end{equation}
на масштабе объединения. В присоединённом представлении три нормированных
индекса также совпадают:
\begin{equation}
 I_3=I_2=I_1^{GUT}=5.
\end{equation}

Критическое различие состоит в том, что gauge action использует не размерность
представления, а квадрат генератора и Dynkin-индекс. Для одного смешанного
блока $X=({\bf3},{\bf2})_{5/6}$ получаем
\begin{equation}
 (I_1,I_2,I_3)=\left(\frac52,\frac32,1\right),
\end{equation}
а для пары $X\oplus\overline X$
\begin{equation}
 (I_1,I_2,I_3)=(5,3,2).
\end{equation}
Следовательно, rank-данные $(6,6/24)$ и gauge-данные $(5/2,3/2,1)$ являются
разными инвариантами.

Инволюция $P$ также не спасает конструкцию. Она различает только чётные блоки
$C\oplus W\oplus Y$ и нечётные $X\oplus\overline X$. Любая функция $f(P)$
оставляет цветовой, слабый и гиперзарядный блоки с одним весом. Добавление
$\operatorname{ad}(Y)^2$ выделяет смешанные блоки, но по-прежнему не различает
$C,W,Y$. Независимые проекторы на эти три сектора возвращают два свободных
относительных gauge-веса после удаления общего масштаба.

\begin{nogocriterion}[Провал rank-action]
Канонический единый trace-action не выводит атласные rank-selector формулы.
Явные секторные проекторы способны их закодировать, но тогда искомый selector
вставляется в действие вручную. Rank-selector сохраняется как экономная
representation-theoretic запись, а не как gauge-action theorem.
\end{nogocriterion}

\section{Gate II: относительный determinant}

В twisted/untwisted отношении локальные heat-kernel части сокращаются. Для
объявленной $\mathbb Z_2$-инволюции изменяются только смешанные блоки
$X\oplus\overline X$; блоки $C,W,Y$ сокращаются между двумя ветвями. Поэтому
любой общий determinant даёт поправку вида
\begin{equation}
 \Delta\boldsymbol{\alpha}^{-1}
 =A_{rel}(5,3,2)
\end{equation}
в порядке $(U(1),SU(2),SU(3))$.

Требуемая текущим gauge-scorecard поправка равна
\begin{equation}
 \Delta\boldsymbol{\alpha}^{-1}_{need}
 =(0.3680,-0.3680,-2.4592).
\end{equation}
Она имеет смешанный знак, тогда как положительный relative determinant лежит
на положительном луче $(5,3,2)$. Наилучшее положительное решение равно нулю и
оставляет $100\%$ исходной нормы ошибки. Даже если искусственно разрешить
произвольный общий знак, наилучшая амплитуда
\begin{equation}
 A_{rel}=-0.11006
\end{equation}
оставляет относительный $L^2$-остаток $96.3\%$.

\begin{proposition}[Direction no-go]
Минимальный $SU(5)$-adjoint relative determinant является конечным и
scheme-safe, но его representation-fixed направление $(5,3,2)$ несовместимо
с требуемой одновременной поправкой электрослабой и сильной связей. Этот вывод
не зависит от общей нормировки и числа одинаковых KK-оболочек.
\end{proposition}

\section{Сохранившийся результат}

Два отрицательных gate не обнуляют проделанную конструкцию:
\begin{itemize}
 \item $SU(5)$-разложение, anomaly-free matter package и $\mathbb Z_2$-чётности
 остаются корректными;
 \item rank-selector остаётся компактным языком спектральных адресов;
 \item relative determinant остаётся чистой топологической response-величиной,
 не зависящей от отдельного локального вычитания.
\end{itemize}

Но ни rank-action, ни минимальный adjoint determinant не являются мостом к
текущим gauge-наблюдаемым. Следующая допустимая попытка должна либо предсказать
новое relative-response наблюдаемое, либо заранее вывести несколько
представлений и их неодинаковые спектры. Использование дополнительных
мультиплетов только ради натяжения threshold-вектора запрещено.

Численные протоколы сохранены в
\texttt{s2t\_su5\_rank\_action\_gate\_results.json} и
\texttt{s2t\_su5\_adjoint\_relative\_determinant\_results.json}.

\end{document}